% !TeX root = ../main.tex

\chapter{数据集、评价指标与方法方案}
\label{ch:method}

\section{数据集介绍}

本文采用 LEVIR-Ship 数据集作为统一实验基准\cite{levirship2021_dataset,chen2022_tinyship_drenet}。该数据集面向中分辨率遥感微小船舶检测任务，能够较好反映复杂海面背景下小目标检出的实际难点。结合当前实验使用的数据组织方式，数据集划分如表~\ref{tab:ch3_dataset_split} 所示。

\begin{table}[htbp]
  \centering
  \caption{LEVIR-Ship 数据集划分}
  \label{tab:ch3_dataset_split}
  \begin{tabular}{lcc}
    \toprule
    数据子集 & 图像数量 & 目标数量 \\
    \midrule
    训练集 & 2321 & 2002 \\
    验证集 & 788 & 665 \\
    测试集 & 788 & 552 \\
    \bottomrule
  \end{tabular}
\end{table}

从任务属性看，LEVIR-Ship 具有目标尺度小、背景纹理复杂、场景差异明显等特点。由于船舶目标在图像中占比很小，仅凭单一指标难以完整描述模型表现，因此本文在实验中同时关注检测精度、召回能力和工程可部署性。

\subsection{数据集样例与任务特征}

从具体场景看，LEVIR-Ship 不仅包含开阔海域中的稀疏目标，也包含近岸区域、港口区域和多目标密集分布场景。不同图像块之间在背景纹理、目标密度、目标朝向和局部对比度上存在明显差异，这使得模型既要具备对弱目标的敏感性，也要具备对复杂背景纹理的抑制能力。对于本文关注的微小船舶检测任务而言，这种“目标小而背景强”的组合恰恰构成了最主要的识别难点。

\begin{figure}[htbp]
  \centering
  \includegraphics[width=0.88\textwidth]{generated/ch3_dataset_examples.png}
  \caption{LEVIR-Ship 数据集典型样例}
  \label{fig:ch3_dataset_examples}
\end{figure}

因此，本文在后续实验中不仅关注主指标 AP50 的变化，还特别重视 Precision、Recall、F1 以及定性样例中的误检和漏检类型。原因在于，对于遥感船舶检测这样兼具检出需求与工程部署需求的任务，仅依赖单一精度指标往往难以完整反映模型是否真正适合系统落地。

\section{数据预处理与统一格式}

为了提升跨框架实验的可比性，本文采用“统一标注协议 + 统一评测口径”的处理思路。首先，对图像与标注文件的一一对应关系进行核验，检查空标注、异常坐标和类别不一致问题。其次，将数据整理为统一的 COCO 中间格式，以便 MMDetection 与 YOLO 系列共享相近的评测流程。再次，固定 train、val、test 划分与样本清单，避免模型间由于划分差异造成额外偏差。最后，将各模型输出统一整理为包含 \texttt{image\_id}、\texttt{model\_name}、\texttt{bbox}、\texttt{score} 和 \texttt{category\_id} 的结构化结果，便于后续可视化展示与数据库存储。

上述处理的核心目的，是尽可能使模型差异主要来源于检测范式本身，而不是数据组织、脚本实现和结果格式差异。对于毕业设计论文而言，这也是保证结论可信度的重要前提。若不同模型分别采用不同的样本划分、不同的评测脚本和不同的结果格式，即便得到若干看似可比较的数字，也很难据此形成可靠结论。因而，本文将“统一格式”视为实验方法的一部分，而不仅仅是训练前的工程准备。

在实现层面，统一格式还有两个直接收益。第一，MMDetection、YOLO 和本地插件推理结果可以通过同一套后处理与可视化逻辑进行整理，降低了跨框架联调成本。第二，结构化结果能够直接复用到 Web 系统数据库中，使实验章中的离线结果与系统章中的在线任务回放使用同一数据表达方式。这一设计保证了论文实验链路与系统实现链路之间的自然衔接。

\section{评价指标}

本文采用 AP50 作为主指标，并辅以 AP50-95、Precision、Recall 和 F1 进行综合评价。相关定义如下。

\begin{equation}
\text{Precision} = \frac{TP}{TP + FP}
\end{equation}

\begin{equation}
\text{Recall} = \frac{TP}{TP + FN}
\end{equation}

\begin{equation}
F1 = \frac{2 \cdot \text{Precision} \cdot \text{Recall}}{\text{Precision} + \text{Recall}}
\end{equation}

\begin{equation}
\text{IoU} = \frac{|B_{pred} \cap B_{gt}|}{|B_{pred} \cup B_{gt}|}
\end{equation}

其中，AP50 表示在 IoU 阈值为 0.5 时的平均精度，主要反映模型在当前任务中的整体检测能力；AP50-95 则在多个 IoU 阈值下进行平均，更强调定位质量与框回归能力。Precision 和 Recall 分别反映误检控制能力与漏检控制能力，F1 用于综合刻画二者的平衡关系。除上述精度类指标外，本文还在后续实验中统计 FPS、参数量和 FLOPs，以辅助分析模型的工程部署属性。

\subsection{指标选择依据}

本文将 AP50 作为主指标，主要是因为微小船舶目标在图像中尺寸较小，当预测框与标注框出现少量偏移时，IoU 会受到较明显影响。若仅使用更严格的高 IoU 阈值进行评价，容易放大框回归误差而弱化“是否检出目标”这一任务核心。对于遥感微小目标检测而言，AP50 更适合首先回答“模型能否稳定发现目标”的问题，而 AP50-95 则用于补充说明“检出的框是否足够准确”。

与此同时，Precision、Recall 与 F1 在本文中具有同等重要的解释价值。原因在于，DRENet、FCOS 和 YOLO26 三类模型在设计目标上并不完全一致，单看 AP50 并不能体现其对误检与漏检的不同取舍。例如，高 Recall 的模型更适合强调漏检控制的场景，而高 Precision 的模型更适合强调输出稳定性与工程部署可控性的场景。因此，本文在后续实验分析中始终把 Precision--Recall 平衡视为结论解释的重要依据，而不是只用单一分数进行简单排序。

\section{模型方案}

为体现不同检测范式在遥感微小船舶任务中的表现差异，本文选择三类代表模型构建对比实验。

\subsection{DRENet（专项方法）}

DRENet 是针对遥感微小船舶任务提出的专项方法，通过退化重建增强机制在训练阶段强化对弱纹理目标的表示能力\cite{chen2022_tinyship_drenet}。该方法对复杂海面背景具有较强针对性，是本文的重要主线模型之一。本文重点关注其在 Recall 和复杂背景场景中的表现。

\subsection{FCOS（Anchor-Free 参考基线）}

FCOS 通过位置回归替代锚框设计，降低了 anchor 尺寸和比例等超参数依赖\cite{tian2019_fcos}。在本文中，FCOS 的角色更接近 anchor-free 参考基线：它用于帮助分析“在不依赖专项结构的情况下，通用 anchor-free 方法在遥感微小船舶任务中能达到怎样的效果”。需要说明的是，当前可用的 FCOS 正式结果采用 MMDetection 默认输入策略，因此其主要承担参考和补充比较作用，而不作为严格同输入尺度公平主结论的核心对象。

\subsection{YOLO26（One-Stage 轻量部署代表）}

YOLO 系列以端到端检测和较高推理效率著称，是工程部署中应用广泛的一类模型\cite{redmon2016_yolo,redmon2018_yolov3,ultralytics2024_docs}。本文采用 YOLO26 作为 one-stage 路线代表，重点考察其在精度、定位质量、训练链路成熟度和系统接入便利性之间的平衡关系。

\section{统一实验协议与实现约束}

为保证实验结论具有可比性，本文对实验协议作了统一约束。第一，使用同一数据划分和同一评测口径。第二，对 DRENet 和 YOLO26 统一采用 512 输入主线；对于 FCOS，则在论文中明确其为历史正式 run 对应的参考基线，并对输入尺度差异进行说明。第三，统一记录实验环境、配置文件、权重路径、日志文件和关键结果，以保证结果可追溯。第四，统一结果输出结构，使算法结果能够直接复用到后续系统实现中。

上述统一协议并不意味着所有模型在每个细节上都做到了完全同构。相反，本文明确保留了模型实现差异的客观边界，并把这些边界写入实验说明。例如，FCOS 当前可用正式结果来自 MMDetection 默认输入策略，因此在论文中只作为 anchor-free 参考基线，而不参与严格同输入尺度的主结论竞争；DRENet 在当前本地插件路径下对非 512 输入存在 shape mismatch，则在消融实验中明确写为实现边界，而不是强行生成不可靠结果。这样的处理方式虽然更“保守”，但更符合本科毕业设计对真实性和可追溯性的要求。

在实现层面，本文通过统一推理入口完成不同模型的封装：

\begin{verbatim}
predict(image_path, model_name, conf_threshold, iou_threshold)
\end{verbatim}

在此基础上，又进一步封装融合推理接口与系统任务执行逻辑。这样做的意义在于，将算法层与系统层解耦，使模型训练、离线评测、单图推理、批量推理和 Web 系统接入共享同一套核心能力，而不需要为不同入口重复编写逻辑。这一设计既有利于代码维护，也使论文中的实验章与系统章形成更自然的衔接。

进一步看，统一推理入口不仅是代码复用问题，也是实验方法问题。若离线评测、命令行可视化与系统任务接口分别维护各自的参数语义和结果结构，就会导致论文结果、演示结果和数据库记录之间出现口径偏移。本文采用统一入口后，可以确保阈值设置、模型选择、结果字段和可视化流程在不同场景下保持一致，从而把“实验结果可信”与“系统展示可信”统一到同一实现基础上。

\section{本章小结}

本章介绍了本文实验所依赖的数据集、统一预处理流程、评价指标和三类模型方案，并说明了统一实验协议与实现约束。通过这些设计，本文建立了面向遥感微小船舶检测任务的可复验实验基础，也为后续的实验结果分析和系统实现提供了统一接口与数据支撑。
